\documentclass[oneside]{ctexbook}
\usepackage{ctex}
\usepackage[utf8]{inputenc}
\usepackage{amsmath}
%\usepackage{amsthm}
\usepackage{ntheorem}
\usepackage{booktabs}
\usepackage{caption}
\usepackage{listings}
\usepackage[dvipsnames]{xcolor}
\usepackage{xeCJK}
\usepackage{bm}
\usepackage{fancyhdr}
\usepackage{graphicx}
\usepackage{amssymb}
\usepackage{mathrsfs}
\usepackage{titlesec, blindtext, color}
\usepackage{arydshln}
\usepackage{hyperref} 
\usepackage[OT1]{fontenc}
\usepackage{geometry}
\usepackage{comment}
\usepackage{extarrows}
\usepackage{inconsolata}
\usepackage{yhmath}
\usepackage{enumitem}
\usepackage[titletoc]{appendix}
\usepackage{verbatim}

\geometry{a4paper,scale=0.8}

\hypersetup{hidelinks,
	colorlinks=true,
	allcolors=black,
	pdfstartview=Fit,
	breaklinks=true
}

\newcounter{sectionCnt} 


\ctexset{
	 chapter = {
		name = {第,章},
		format = \linespread{1.0}\zihao{1}\bfseries\sffamily\centering,
		nameformat = {\fontsize{30pt}{0pt}\centering\vspace{15pt}},
		titleformat = {\centering},
		number = \chinese{chapter},
		numberformat = \sffamily,
		aftername = ,
		beforeskip = {7pt},
		afterskip = {18pt},
		aftername = {\\},
		pagestyle = plain,
	},
	section={
		%format用于设置章节标题全局格式，作用域为标题和编号
		%字号为小三，字体为黑体，左对齐
		%+号表示在原有格式下附加格式命令
		format+ = \zihao{2} \heiti \raggedright\centering,
		%name用于设置章节编号前后的词语
		%前、后词语用英文状态下,分开
		%如果没有前或后词语可以不填
		name = {,.\,},
		%number用于设置章节编号数字输出格式
		%\chinese输出section编号为中文，\arabic输出为阿拉伯数字
		%number = \Roman{sectionCnt},
		number = \chinese{sectionCnt},
		%beforeskip用于设置章节标题前的垂直间距
		%ex为当前字号下字母x的高度
		%基础高度为1.0ex，可以伸展到1.2ex，也可以收缩到0.8ex
		beforeskip = 1.0ex plus 0.2ex minus .2ex,
		%afterskip用于设置章节标题后的垂直间距
		afterskip = 1.0ex plus 0.2ex minus .2ex,
		%aftername用于控制编号和标题之间的格式
		%\hspace用于增加水平间距
		aftername = \hspace{0pt}
	},
	subsection={
		format+ = \zihao{-2} \kaishu  \raggedright,
		%仅输出subsection编号且为中文
		number = \roman{subsection},
		name = {,.},
		beforeskip = 1.0ex plus 0.2ex minus .2ex,
		afterskip = 1.0ex plus 0.2ex minus .2ex,
		aftername = \hspace{0pt}
	},
	subsubsection={
		%设置对齐方式为居中对齐
		format+ = \zihao{-3} \fangsong,
		%仅输出subsubsection编号，格式为阿拉伯数字，打字机字体
		number = \ttfamily\arabic{subsubsection},
		name = {,.},
		%beforeskip = 1.0ex plus 0.2ex minus .2ex,
		afterskip = 1.0ex plus 0.2ex minus .2ex,
		aftername = \hspace{0pt}
	}
}

\newenvironment{Appendices}[1]
{

	\ctexset{
		chapter = {
			name = {#1},
			format = \linespread{1.0}\zihao{1}\bfseries\sffamily\centering,
			nameformat = {\fontsize{30pt}{0pt}\centering\vspace{15pt}},
			titleformat = {\centering},
			%number = \chinese{chapter},
			number = ,
			numberformat = \sffamily,
			aftername = ,
			beforeskip = {7pt},
			afterskip = {18pt},
			aftername = {\\},
			pagestyle = plain,
		},
		section={
			%format用于设置章节标题全局格式，作用域为标题和编号
			%字号为小三，字体为黑体，左对齐
			%+号表示在原有格式下附加格式命令
			format+ = \zihao{-1} \heiti \raggedright\centering,
			%name用于设置章节编号前后的词语
			%前、后词语用英文状态下,分开
			%如果没有前或后词语可以不填
			name = {,.\,},
			%number用于设置章节编号数字输出格式
			%\chinese输出section编号为中文，\arabic输出为阿拉伯数字
			number = \Roman{section},
			%beforeskip用于设置章节标题前的垂直间距
			%ex为当前字号下字母x的高度
			%基础高度为1.0ex，可以伸展到1.2ex，也可以收缩到0.8ex
			beforeskip = 1.0ex plus 0.2ex minus .2ex,
			%afterskip用于设置章节标题后的垂直间距
			afterskip = 1.0ex plus 0.2ex minus .2ex,
			%aftername用于控制编号和标题之间的格式
			%\hspace用于增加水平间距
			aftername = \hspace{0pt}
		},
		subsection={
			format+ = \zihao{-2} \kaishu  \raggedright,
			%仅输出subsection编号且为中文
			number = \roman{subsection},
			name = {,.},
			beforeskip = 1.0ex plus 0.2ex minus .2ex,
			afterskip = 1.0ex plus 0.2ex minus .2ex,
			aftername = \hspace{0pt}
		},
		subsubsection={
			%设置对齐方式为居中对齐
			format+ = \zihao{-3} \fangsong,
			%仅输出subsubsection编号，格式为阿拉伯数字，打字机字体
			number = \ttfamily\arabic{subsubsection},
			name = {,.},
			%beforeskip = 1.0ex plus 0.2ex minus .2ex,
			afterskip = 1.0ex plus 0.2ex minus .2ex,
			aftername = \hspace{0pt}
		}
	}
}

\fancyhf{}
\fancyhead[L]{
	\begin{minipage}[c]{0.06\textwidth}
		\hyperlink{Index}{\includegraphics[height=7.5mm]{1.jpg}}
	\end{minipage}
	\begin{minipage}[c]{0.4\textwidth}
	线代冲刺笔记
	\end{minipage}}
\fancyhead[R]{
    \begin{minipage}[r]{0.1\textwidth}
	\href{https://qifengggg.github.io/}{奇峰}
    \end{minipage}
}
\cfoot{\hyperlink{Index}{\thepage}}

\pagestyle{fancy}

\fancypagestyle{plain}{
  \fancyhead{}
  \renewcommand{\headrulewidth}{0pt}
}
%\renewcommand{\qedsymbol}{$\blacksquare$}

\renewcommand{\labelenumi}{\roman{enumi}.}
\renewcommand{\labelitemii}{$ \circ $ }
\newcommand{\dis}{\displaystyle}
\newcommand{\Attention}[1]{\textcolor{red}{\bf{#1}}}
\newcommand{\red}{\textcolor{red}}
\newcommand{\sssubsection}[1]{\noindent\textbf{#1}}
\newcommand{\nextline}{\vspace{5pt}\newline}
\newcommand{\dsqrt}{\dis\sqrt}
\newcommand{\pprime}{{\prime\prime}}
\newcommand{\whichis}{\mathop=\limits^\triangle}
\newcommand{\R}{\mathbb R}
\newcommand{\qline}{\underline{~~~~~~~~~~}}

\newcommand{\Section}{\stepcounter{sectionCnt}\section}

\title{\includegraphics[scale=0.6]{1.jpg}\\ \textsf{线代冲刺笔记}}
\author{\href{https://qifengggg.github.io/}{奇峰}}
\date{之前}

\begin{document}

\theoremseparator{}
\newtheorem{def1}{定义}[section]
\newtheorem{theo1}{定理}[section]
\newtheorem{func1}{方法}[section]
\newtheorem{infer1}{推论}[section]
\newenvironment{proof}{\begin{itemize}\item \textbf{证明}
	
	}{\end{itemize}}
%\newenvironment{<环境名称>}[<参数个数>][<首参数默认值>]{<环境前定义>}
%							{<环境后定义>}


\newenvironment{Def}[1][\quad{}]{\begin{def1}\textbf{#1}}{\end{def1}}
\newenvironment{Theo}[1][\quad{}]{\begin{theo1}\textbf{#1}}{\end{theo1}}
\newenvironment{Func}[1][\quad{}]{\begin{func1}\textbf{#1}}{\end{func1}}
\newenvironment{Infer}[1][\quad{}]{\begin{infer1}\textbf{#1}}{\end{infer1}}
\newenvironment{Field}[1][\quad{}]{\noindent\newline\textbf{#1}}{}

\setlength{\parskip}{3pt}
\setenumerate[1]{topsep=2pt}
\setitemize[1]{topsep=2pt}

\frontmatter
\maketitle
\hypertarget{Index}{}
\tableofcontents

\mainmatter

\chapter{二重积分}

\section{交换积分次序}

$ \begin{cases}
    \textrm{直角 - 穿针法}\\
    \textrm{极坐标}\begin{cases}
        xy\textrm{坐标系 - 作圆弧}\\ 
        \theta r\textrm{坐标系 - 穿针法(当成直角坐标)}
    \end{cases}
\end{cases} $ 

二重变限积分求极限时，
$ \begin{cases}
    \textrm{洛必达结合交换积分次序}\\\textrm{二重积分中值定理(非零因子)}
\end{cases} $ 

\sssubsection{例题T6.2.2}

\begin{equation*}
    \begin{aligned}
        {\displaystyle\lim_{t\rightarrow 0^+}}\dfrac{1}{t^3}\int_0^t\mathrm{d}x
        \int_{x^2}^{t^2}\arctan(\cos(2x+3\dsqrt{y}))\mathrm{d}y &= 
        {\displaystyle\lim_{t\rightarrow 0^+}}\dfrac{1}{t^3}\iint\limits_{D}
        \arctan(\cos(2x+3\dsqrt{y}))\mathrm{d}x\mathrm{d}y
        \\&= {\displaystyle\lim_{t\rightarrow 0^+}}\dfrac{1}{t^3}
        \arctan(\cos(2\xi+3\dsqrt{\eta}))\cdot S_D \,(\xi,\eta \in D)
        \\&= {\displaystyle\lim_{t\rightarrow 0^+}}\dfrac{1}{t^3}
        \arctan(\cos(2\xi+3\dsqrt{\eta}))\cdot\dfrac{2}{3}t^3
        \\&= \dfrac{\pi}{4}\cdot\dfrac{2}{3} = \dfrac{\pi}{6}.
    \end{aligned}
\end{equation*}

\sssubsection{例题T6.3.2}

\begin{equation*}
    \begin{aligned}
        \int_0^\frac{\pi}{2}\dfrac{1}{\sqrt{x}}\mathrm{d}x\int_{\sqrt{x}}^{\sqrt{\frac{\pi}{2}}}
        \dfrac{1}{1+(\tan y^2)^{\sqrt 2}}\mathrm{d}y 
        &= \int_0^{\sqrt{\frac{\pi}{2}}}\dfrac{1}{1+(\tan y^2)^{\sqrt 2}}
        \mathrm{d}y\int_0^{y^2}\dfrac{1}{\sqrt x}\mathrm{d}x
        \\&= \int_0^{\sqrt{\frac{\pi}{2}}}\dfrac{2y}{1+(\tan y^2)^{\sqrt 2}}\mathrm{d}y
        \\&= \int_0^{\frac{\pi}{2}}\dfrac{\mathrm{d}t}{1+(\tan t)^{\sqrt 2}}
        \\&= \int_0^{\frac{\pi}{2}}\dfrac{(\cos t)^{\sqrt 2}\mathrm{d}t}{(\sin t)^{\sqrt 2}+(\cos t)^{\sqrt 2}} 
        \\\textrm{区间再现}&=\dfrac{1}{2}\cdot\dfrac{\pi}{2} = \dfrac{\pi}{4}.
    \end{aligned}
\end{equation*}

\sssubsection{例题T6.3.4}

\begin{equation*}
    \begin{aligned}
        \int_0^1\dfrac{x^b-x^a}{\ln x}\mathrm{d}x &= 
        \int_0^1 \mathrm{d}x\int_a^b x^y\mathrm{d}y \\ 
        &= \int_a^b \mathrm{d}y \int_0^1 x^y\mathrm{d}x
        \\&= \int_a^b \dfrac{\mathrm{d}y}{1+y} = \ln\dfrac{1+b}{1+a}.
    \end{aligned}
\end{equation*}

\sssubsection{例题T6.3.5}

此处以 $ \theta - r $ 坐标系交换积分次序。

\begin{equation*}
    \begin{aligned}
        \int_0^{2\pi}(\theta^2-1)\mathrm{d}\theta\int_{\frac{\theta}{2}}^\pi e^{r^2}\mathrm{d}r &=
        \int_0^\pi e^{r^2}\mathrm{d}r\int_0^{2r}(\theta^2-1)\mathrm{d}\theta \\ 
        &= \dfrac{8}{3}\int_0^\pi r^3e^{r^2}\mathrm{d}r - 2\int_0^\pi re^{r^2}\mathrm{d}r
        \\ &= \dfrac{1}{3}(4\pi^2-7)e^{\pi^2}+\dfrac{7}{3}.
    \end{aligned}
\end{equation*}

\section{二重积分的计算}

考虑
$ \begin{cases}
    \textrm{直角坐标系}\\
    \textrm{极坐标变换}\\
    \textrm{奇偶性}\\
    \textrm{轮换对称性}\\
    \textrm{形心公式}\\
    \textrm{平移变换}\\
    \textrm{雅可比行列式}
\end{cases} $ 

\sssubsection{被积函数为齐次函数 $ f(\dfrac{y}{x}) $ T6.6}

\begin{enumerate}
    \item[\textbf{例题}] 设有区域 $ D = \{(x,y)|x|+|y|\leq 1\}, $ 
    计算 $ \dis \iint\limits_D e^{\dfrac{|y|}{|x|+|y|}}\mathrm{d}x\mathrm{d}y. $ 
    \item[\textrm{\textbf{方法}}] 可以知道，
    \begin{equation*}
        \begin{aligned}
            \textrm{原积分} &= 4\iint\limits_{D_1}e^{\dfrac{y}{x+y}}\mathrm{d}x\mathrm{d}y \\ 
            &= 4\int_0^{\frac{\pi}{2}}\mathrm{d}\theta\int_0^{\frac{1}{\sin \theta + \cos \theta}}
            e^{\frac{\sin \theta}{\sin \theta + \cos \theta}}r\mathrm{d}r
            \\&= 2\int_0^{\frac{\pi}{2}} e^{\frac{\sin \theta}{\sin \theta + \cos \theta}} 
            \dfrac{1}{(\cos \theta + \sin \theta)^2}\mathrm{d}\theta
            \\&= 2\int_0^{\frac{\pi}{2}} e^{\frac{\sin \theta}{\sin \theta + \cos \theta}} 
            \mathrm{d}\dfrac{\sin \theta}{\sin \theta + \cos \theta}
            \\&= 2e^{\frac{\sin \theta}{\sin \theta + \cos \theta}}
            \Big|_0^\frac{\pi}{2} = 2(e-1).
        \end{aligned}
    \end{equation*}
\end{enumerate}

事实上，对于包含齐次函数 $ f(\dfrac{y}{x}) $ 的二重积分，都应该考虑利用极坐标变出 
$ \mathrm{d}\dfrac{A\sin\theta+B\cos\theta}{\sin\theta+\cos\theta}. $

对于 $ f(\dfrac{y^2}{x^2}),f(\dfrac{ax+by}{x+y}) $ 也是类似。

\sssubsection{极坐标压轴题 - T6.7}

\begin{enumerate}
    \item[\textbf{例题}] 设 $ f(x,y) $ 在区域 $ D = \{(x,y)|x^2+y^2\leq 1\} $ 
    上有一阶连续偏导数，在 $ D $ 的边界上 $ f(x,y) $ 恒为零，求
    $$
        {\displaystyle\lim_{\varepsilon\rightarrow 0^+}}\iint\limits_{\varepsilon^2\leq x^2+y^2\leq 1}
        \dfrac{xf'_x+yf'_y}{x^2+y^2}\mathrm{d}x\mathrm{d}y.
    $$
    \item[\textbf{方法}] 引入极坐标，则有
    \begin{equation*}
        \begin{aligned}
            &\dfrac{\partial f}{\partial \rho} = 
            \dfrac{\partial f}{\partial x}\dfrac{\partial x}{\partial \rho}
            + \dfrac{\partial f}{\partial y}\dfrac{\partial y}{\partial \rho}
            = f'_x\cos\theta + f'_y\sin\theta \\ 
            &\dfrac{\partial f}{\partial \theta} = 
            \dfrac{\partial f}{\partial x}\dfrac{\partial x}{\partial \theta}
            + \dfrac{\partial f}{\partial y}\dfrac{\partial y}{\partial \theta}
            = -f'_x(\rho \sin\theta)+f'_y(\rho \cos\theta)\\ 
            \Rightarrow & \rho \dfrac{\partial f}{\partial \rho} = 
            \rho\left[f'_x\cos\theta + f'_y\sin\theta\right]
            =xf'_x+yf'_y.
        \end{aligned}
    \end{equation*}
    那么，对于题设，有
    \begin{equation*}
        \begin{aligned}
            \textrm{原极限}&= {\displaystyle\lim_{\varepsilon\rightarrow 0^+}}
            \int_0^{2\pi}\mathrm{d}\theta\int_{\varepsilon}^1
            \dfrac{\partial f}{\partial \rho}\mathrm{d}\rho \\ 
            &= {\displaystyle\lim_{\varepsilon\rightarrow 0^+}}
            \int_0^{2\pi}f(\rho\cos\theta,\rho\sin\theta)\Big|_\varepsilon^1\mathrm{d}\theta\\
            &= -{\displaystyle\lim_{\varepsilon\rightarrow 0^+}}
            \int_0^{2\pi}f(\varepsilon\cos\theta,\varepsilon\sin\theta)\mathrm{d}\theta \\ 
            &= -\int_0^{2\pi}f(0,0)\mathrm{d}\theta = -2\pi f(0,0).
        \end{aligned}
    \end{equation*}
\end{enumerate}

\sssubsection{极坐标逆问题 - T6.8}

在极坐标不好做的情况下，将 $ (\rho,\theta) $ 转换回 $ (x,y) $ 即可。

\newpage

\sssubsection{奇偶性 - T6.9-10}

\begin{itemize}
    \item \textbf{奇偶性的推广}
    
    设 $ f(x,y) $ 在区域 $ D = D_1\bigcup D_2 $ 上连续，其中 $ D_1,D_2 $ 关于原点对称，则
    $$
        \iint\limits_Df(x,y)\mathrm{d}x\mathrm{d}y
        =\begin{cases}
            2\dis\iint\limits_{D_1}f(x,y)\mathrm{d}x\mathrm{d}y,\quad{} f(-x,-y) = f(x,y)\\
            0,\quad{} f(-x,-y) = -f(x,y)\\
        \end{cases}
    $$
\end{itemize}

\begin{enumerate}
    \item[\textbf{例题}]
    设区域 $ D $ 由曲线 $ y = x^3,y = x $ 围成。求
    $$
        \iint\limits_D\left[e^{x^2}+\sin(x+y)\right]\mathrm{d}x\mathrm{d}y.
    $$
    \item[\textbf{方法}] 
    由奇偶性，可以知道
    \begin{equation*}
        \begin{aligned}
            \textrm{原积分}&=2\int_0^1\mathrm{d}x\int_{x^3}^x e^{x^2}\mathrm{d}y\\ 
            &= 2\int_0^1e^{x^2}(x-x^3)\mathrm{d}x 
            \\&= \int_0^1(1-x^2)e^{x^2}\mathrm{d}x^2 = e-2.
        \end{aligned}
    \end{equation*}
\end{enumerate}

\begin{enumerate}
    \item[\textbf{例题}]
    对 $ D = \{(x,y)|x^2+y^2\leq 1, x+y\geq 0\}, $ 求
    $$
        \iint\limits_D\dfrac{1+2x^2+xy^2}{1+x^2+y^2}\mathrm{d}x\mathrm{d}y.
    $$
    \item[\textbf{方法}] 
    原积分等于 $ \dis \iint\limits_D\dfrac{1+2x^2}{1+x^2+y^2}\mathrm{d}x\mathrm{d}y+
    \iint\limits_D\dfrac{xy^2}{1+x^2+y^2}\mathrm{d}x\mathrm{d}y. $ 
    前一部分利用轮换对称性。对后一部分，以 $ y = x $ 将 $ D $ 
    划分为 $ D_1,D_2, $ 以 $ y = 0 $ 将 $ D_1 $ 划分为 $ D_{11},D_{12}. $ 
    那么，
    \begin{equation*}
        \begin{aligned}
            \textrm{后半部分积分} &= \iint\limits_{D_1}+\iint\limits_{D_2}
            \dfrac{xy^2}{1+x^2+y^2}\mathrm{d}x\mathrm{d}y\\ 
            &= 2\iint\limits_{D_12}\dfrac{xy^2}{1+x^2+y^2}\mathrm{d}x\mathrm{d}y\\ 
            &= 2\int_{0}^\frac{\pi}{4}\mathrm{d}\theta\int_0^1
            \dfrac{r^4\sin^2\theta\cos\theta}{1+r^2}\mathrm{d}r = \dfrac{1}{3\sqrt 2}
            \left(\dfrac{\pi}{4}-\dfrac{2}{3}\right).
        \end{aligned}
    \end{equation*}
\end{enumerate}

\sssubsection{轮换对称性的推广 - T6.11}

\begin{itemize}
    \item \textbf{轮换对称性的推广}
    
    $$
        \iint\limits_Df(x,y)\mathrm{d}\sigma = 
        \iint\limits_{D'}f(y,x)\mathrm{d}\sigma
    $$
    其中 $ D' = \{(y,x)|(x,y)\in D\}. $ 
\end{itemize}
\chapter{无穷级数}

\section{数项级数}

\subsection{数项级数概念}

\begin{Def}[数项级数]

    无穷多个数 $ u_1,u_2,\dots,u_n,\dots $ 按照一定的顺序相加得到的表达式
     $ \sum_{n=1}^\infty u_n $ 称为数项级数，简称级数，其中 $ u_n $ 称为通项。
\end{Def}

级数是无数个的，也是有序的。

\begin{Def}[级数收敛]

    称 $ S_n = \sum_{k=1}^n u_k $ 为级数的部分和，数列 $ \{S_n\} $ 即为部分和数列。

    若 $ {\displaystyle\lim_{n\rightarrow \infty}}S_n\xlongequal{\exists}S $ ，
    则称级数 $\dis \sum_{n=1}^\infty u_n $ 收敛，且和为 $ S $ ;反之，若极限不存在，则称其
    为发散的，发散级数没有和的概念。
\end{Def}

以下是几个重要级数及其敛散性。

\begin{itemize}
    \item 几何级数 - $\dis \sum_{n=0}^\infty ar^n $ ，当 $ |r|<1 $ 时，其收敛于 $ \dfrac{a}{1-r} $ ;
    若 $ |r|\geq1 $ ，则其发散。
    \item p级数 - $\dis \sum_{n=1}^\infty \dfrac{1}{n^p} $ ，若 $ p > 1 $ ，其收敛；若 $ p\leq 1 $ ，其发散。
\end{itemize}

事实上，有 $\dis \sum_{n=1}^\infty \dfrac{1}{n^2} = \dfrac{\pi^2}{6}. $ 

\section{级数基本性质及其收敛必要条件}

级数有以下性质。

\begin{itemize}
    \item 若 $ \dis \sum_{n=1}^{\infty}u_n,\sum_{n=1}^{\infty}v_n $ 都收敛，则
    $ \dis \sum_{n = 1}^{\infty}(u_n\pm v_n) $ 也收敛，且有
    $ \dis \sum_{n = 1}^{\infty}(u_n\pm v_n) = \sum_{n=1}^{\infty}u_n\pm\sum_{n=1}^{\infty}v_n $ .
    \item 在级数中增加、减少、改变有限项不会影响级数敛散性。
    \item 收敛级数具有结合律，即可以对级数的项任意地添加括号，新级数仍然收敛于同一个和。
    注意，去括号不成立。
    \item 对 $ k \neq 0 $ 有 $ \dis \sum_{n=1}^{\infty}ku_n $ 与  $ \dis \sum_{n=1}^{\infty}u_n $
    具有相同的敛散性。
    \item 级数  $ \dis \sum_{n=1}^{\infty}u_n $ 收敛的必要条件是 $ {\displaystyle\lim_{n\rightarrow \infty}}u_n = 0 $ .
\end{itemize}

\section{正项级数及其敛散性判别法}

\begin{Def}[正项级数]

    若 $ u_n\geq 0 $ ，则有 $\dis \sum_{n=1}^{\infty}u_n $ 称为正项级数。
\end{Def}

\begin{Theo}[正项级数收敛基本定理]

    正项级数 $\dis \sum_{n=1}^{\infty}u_n $ 收敛 $ \Leftrightarrow  S_n $ 有上界。
\end{Theo}

正项级数有以下几种判别法。

\begin{Theo}[直接比较法]

    设当 $ n\geq N $ 时，有 $ 0\leq u_n\leq v_n $ 总是成立，则\begin{itemize}
        \item 若 $\dis \sum_{n=1}^{\infty}v_n $ 收敛，则 $\dis \sum_{n=1}^{\infty}u_n $ 收敛；
        \item 若 $\dis \sum_{n=1}^{\infty}u_n $ 发散，则 $\dis \sum_{n=1}^{\infty}v_n $ 发散。
    \end{itemize}
\end{Theo}

\begin{Theo}[间接比较法]

    设 $ u_n\geq 0, v_n>0,n = 1,2,3,\dots $，且有 $ {\displaystyle\lim_{n\rightarrow \infty}}
    \dfrac{u_n}{v_n} = A  $ ，则\begin{itemize}
        \item 若 $ 0<A<+\infty $ ，两级数具有相同的敛散性；
        \item 若 $ A = 0 $ ，则若 $\dis \sum_{n=1}^{\infty}v_n $ 收敛，有 $\dis \sum_{n=1}^{\infty}u_n $ 收敛；
        \item 若 $ A = +\infty $ ，若 $\dis \sum_{n=1}^{\infty}u_n $ 收敛，则 $\dis \sum_{n=1}^{\infty}v_n $ 收敛。
    \end{itemize}
\end{Theo}

\begin{Theo}[比值判别法]

    设 $ u_n > 0 $ ，有 $ {\displaystyle\lim_{n\rightarrow \infty}}\dfrac{u_{n+1}}{u_{n}} = \rho $ ，
    则\begin{itemize}
        \item 若 $ \rho > 1 $ ，有 $\dis \sum_{n=1}^{\infty}u_n $ 收敛；
        \item 若 $ \rho < 1 $ ，则其发散；
        \item 若 $ \rho = 1 $ ，该判别法失效。
    \end{itemize}
\end{Theo}

\begin{Theo}[根值判别法]

    设 $ u_n \leq 0 $ ，有 $ {\displaystyle\lim_{n\rightarrow \infty}}\sqrt[n]{u_n} = \rho $ ，
    则\begin{itemize}
        \item 若 $ \rho > 1 $ 甚至为$ +\infty $ 时，有 $\dis \sum_{n=1}^{\infty}u_n $ 收敛；
        \item 若 $ \rho < 1 $ ，则其发散；
        \item 若 $ \rho = 1 $ ，该判别法失效。
    \end{itemize}
\end{Theo}

\begin{Theo}[积分准则]

    若函数在 $ [1,+\infty) $ 上非负、连续、递减，则级数 $\dis \sum_{n=1}^{\infty}f(n) $ 与
    积分 $ \dis \int_1^{+\infty}f(x)\mathrm{d}x $ 具有相同的敛散性。
\end{Theo}

注意，判别正项级数敛散性时，\begin{equation*}
    \begin{matrix}
        \textrm{间接比较法} \\\Downarrow\\ 
        \textrm{比值法}\\\Downarrow\\
         \textrm{根值法}\\\Downarrow\\
          \textrm{积分准则}\\\Downarrow\\
        \textrm{直接比较法}\\\Downarrow\\
         \textrm{定义·部分和数列}\\\Downarrow\\
          \textrm{收敛级数定义}
    \end{matrix}
\end{equation*}

此外，若通项为积分或者是抽象的，从直接比较法开始应用，试图放缩被积函数以去除积分号。

\section{任意项级数}

\subsection{交错级数及其敛散性判别法}

\begin{Def}[交错级数]

    若 $\dis u_n>0,\sum_{n=1}^{\infty}(-1)^{n+1}u_n $ 称为交错级数。
\end{Def}

\begin{Theo}[莱布尼茨判别法]

    交错级数 $\dis \sum_{n=1}^{\infty}(-1)^{n+1}u_n $ 满足
    \begin{itemize}
        \item $ u_{n+1}\leq u_n $ 
        \item $ {\displaystyle\lim_{n\rightarrow \infty}}u_n = 0 $ 
    \end{itemize}
    则该交错级数收敛，且 $\dis 0\leq \sum_{n=1}^{\infty}(-1)^{n+1}u_n\leq u_1 $ .
\end{Theo}

\subsection{条件收敛与绝对收敛}

\begin{Def}[条件收敛与绝对收敛]

    若 $ \dis \sum_{n=1}^{\infty}|u_n| $ 收敛，则称 $ \dis \sum_{n=1}^{\infty}u_n $ 绝对收敛。

    若 $ \dis \sum_{n=1}^{\infty}u_n $ 收敛而 $ \dis \sum_{n=1}^{\infty}|u_n| $ 发散，则
    $ \dis \sum_{n=1}^{\infty}u_n $ 称为条件收敛。
\end{Def}

绝对收敛的级数有交换律，而条件收敛的级数没有。

\begin{Theo}[条件收敛必收敛]

    $ \dis \sum_{n=1}^{\infty}|u_n|$ 收敛 $\dis \Rightarrow \sum_{n=1}^{\infty}u_n $ 收敛。
\end{Theo}

\section{幂级数}

\subsection{函数项级数}

\begin{Def}[函数项级数]

    设 $ u_n(x),n=1,2,\dots $ 定义在 $ I $ 上，则 $ \dis \sum_{n=1}^{\infty}u_n(x) $ 
    称为区间 $ I $ 上的函数项级数。
\end{Def}

\begin{Def}[收敛点与收敛域]

    对 $ x_0\in I $ ，若常数项级数 $\dis \sum_{n=1}^{\infty}u_n(x_0) $ 收敛，则称 $ x_0 $ 为函数项级数
    $ \dis \sum_{n=1}^{\infty}u_n(x) $ 的收敛点；反之，若其发散，则为发散点。

    函数项级数 $ \dis \sum_{n=1}^{\infty}u_n(x) $ 的全部收敛点构成的集合是收敛域，全部发散点构成的集合
    是发散域。
\end{Def}

\begin{Def}[和函数]

    和函数 $\dis S(x) =: \sum_{n=1}^{\infty}u_n(x) $ ，其定义域为收敛域。
\end{Def}

\subsection{幂级数}

\begin{Def}[幂级数]

    形如$ \dis \sum_{n=1}^\infty a_n(x-x_0)^n $ 的级数称为 $ (x-x_0) $ 的幂级数，常数 $ a_n $ 称为
    幂级数的系数。
\end{Def}

显然，当 $ x_0 = 0 $ 时，有 $ \dis \sum_{n=0}^{\infty}a_nx^n $ 称为 $ x $ 的幂级数。

\begin{Theo}[阿贝尔定理]

    若幂级数 $\dis \sum u_n(x) $ 在 $ x = x_0 $ 处收敛，则对任意 $ x $ 满足 $ |x| < |x_0| $ ,
    都有$ \dis \sum u_n(x) $ 绝对收敛；
    若在 $ x = x_0 $ 发散，则对任意 $ x $ 满足 $ |x|>|x_0| $ 都有$\dis \sum u_n(x) $ 发散。
\end{Theo}

\subsection{幂级数收敛半径及其求法}

幂级数的收敛半径 $ R $ 可以由以下公式求出。

\begin{itemize}
    \item $ \dis R = {\displaystyle\lim_{n\rightarrow \infty}}\left|\dfrac{a_n}{a_{n+1}}\right| $ 
    \item $ \dis R = {\displaystyle\lim_{n\rightarrow \infty}}\dfrac{1}{\dsqrt[n]{|a_n|}} $ 
\end{itemize}

注意，公式仅对标准幂级数成立，即要求幂函数的次数逐次增长；

收敛区间的中心为 $ x_0 $. 以 $ x = 0 $ 为例子，若存在收敛半径 $ R\geq 0 $ ，则
$ x\in (-R,R) $ 时级数绝对收敛，$ x\in R\big\backslash(-R,R) $ 时级数发散。
若级数有条件收敛点，则其只能在 $ \pm R $ 处，因而最多有两个。

在求收敛半径后，必须写出收敛区间，然后判断端点处敛散性。

\subsection{幂级数的性质}

\sssubsection{幂级数的加减运算}

设 $ \dis \sum_{n=1}^\infty a_nx^n = f(x),|x|< R_1;\sum_{n=1}^\infty b_nx^n = g(x), |x|< R_2 $，则有
$ \dis \sum_{n=1}^\infty (a_n\pm b_n)x^n = f(x)\pm g(x), |x|<\min(R_1,R_2)  $ .

\sssubsection{幂级数的逐项求导与逐项积分}

设幂级数 $ \dis \sum_{n=0}^\infty a_nx^n $ 的收敛半径 $ R>0 $ 且有和函数 $ \dis S(x) = \sum_{n=1}^\infty a_nx^n $ ，则
成立以下性质。
\begin{itemize}
    \item $ S(x) $ 在 $ \overset{\textrm{收敛区间}}{(-R,R)} $ 内可导，且有逐项求导公式$$
        S'(x) = \left(\sum_{n=0}^\infty a_nx^n\right)' = \sum_{n=0}^\infty (a_nx^n)' = \sum_{n=0}^\infty na_nx^{n-1}
    $$ 
    求导后幂级数收敛半径不变，因而 $ S(x) $ 在 $ (-R,R) $ 内有任意阶导数；
    \item $ S(x) $ 在$ \overset{\textrm{收敛区间}}{(-R,R)} $内有逐项积分公式$$
        \int_0^x S(t)\mathrm{d}t = \sum_{n=0}^\infty\int_0^x a_nt^n\mathrm{d}t = \sum_{n=0}^\infty \dfrac{a_n}{n+1}x^{n+1}
    $$ 
    积分后收敛半径不变；
    \item 幂级数的和函数 $ S(x) $ 在收敛域上连续。
\end{itemize}

\sssubsection{单位幂级数}

$ \dis \sum_{n=0}^\infty x^n $ 称为单位幂级数,显然其收敛半径是1,且收敛区间内，有 $ \dis \sum x^n = \dfrac{1}{1-x} $ .

$ \dis \sum_{n=0}^\infty (-1)^{n+1}x^n $ 收敛半径也是1, 且收敛区间内，有 $ \dis \sum x^n = \dfrac{1}{1+x} $ .


\section{函数展开成幂级数}

\subsection{概念}

设 $ f(x) $ 是给定的函数，若存在幂级数 $ \dis \sum_{n=0}^\infty a_n(x-x_0)^n $ 使得任意 $ x\in I $ 
级数 $ \dis \sum_{n=0}^\infty a_n(x-x_0)^n $ 都收敛到 $ f(x) $ ，即 $\dis \forall x\in I, f(x) = \sum_{n=0}^\infty a_n(x-x_0)^n $ ,
则称 $ f(x) $ 在区间 $ I $ 上可展开为 $ x - x_0 $ 的幂级数，称级数 $ \dis \sum_{n=0}^\infty a_n(x-x_0)^n $ 为$ f(x) $ 的幂级数。

\begin{Def}[幂级数展开唯一性]

    若 $ f(x) $ 在 $ x = x_0 $ 可以展开为幂级数 $ \dis \sum_{n=0}^\infty a_n(x-x_0)^n $ ，则该展开式是唯一的，且为
    泰勒级数，即有 $ \dis f(x) = \sum_{n=0}^\infty a_n(x-x_0)^n = \sum_{n=0}^\infty \dfrac{f^{(n)(x_0)}}{n!}(x-x_0)^n$
\end{Def}

\subsection{函数展开成幂级数的条件及其形式}

\sssubsection{泰勒级数与马克劳林级数的定义}

设函数 $ f(x) $ 在 $ x_0 $ 的某一领域 $ |x-x_0|<\sigma $ 内具有任意阶导数，则级数
$ \dis \sum_{n=0}^\infty \dfrac{f^{(n)(x_0)}}{n!}(x-x_0)^n $ 称为函数 $ f(x) $ 在 $ x = x_0 $ 处的泰勒级数。

特别地，若 $ x_0 = 0 $ ，称级数为$ f(x) $ 的马克劳林级数。

\sssubsection{函数展开成为幂级数的充要条件}

设 $ f(x) $ 在 $ |x-x_0|<R $ 内有任意阶导数，其泰勒公式为$$
    f(x) = \sum_{i=0}^n \dfrac{f^{(i)}(x_0)}{i!}(x-x_0)^i + R_n(x)
$$ 其中$ R_n(x) $为 $ n $ 阶余项，
其拉格朗日型为 $ \dis R_n(x) = \dfrac{f^{(n+1)}[x_0+\theta(x-x_0)]}{(n+1)!}(x-x_0)^{n+1},\theta\in(0,1) $ ,
则 $ \dis f(x) = \sum_{n=0}^\infty \dfrac{f^{(n)}(x_0)}{ni!}(x-x_0)^n,|x-x_0|<R $ 的充要条件为
$ {\displaystyle\lim_{n\rightarrow \infty}}R_n(x) = 0,|x-x_0|<R $ .

特别地，当 $ x_0 = 0 $ 时得到函数展开为马克劳林级数的充要条件。

\subsection{函数展开成为幂级数的方法}

\begin{itemize}
    \item 公式法 $$
     \dis f(x) = \sum_{n=0}^\infty \dfrac{f^{(n)}(x_0)}{ni!}(x-x_0)^n,|x-x_0|<R 
    $$
    \item 变量替换；
    \item 逐项求导；
    \item 逐项积分；
    \item 恒等变形、变量代换等。
\end{itemize}

幂级数有8个常用的展开式。

\begin{itemize}
    \item $ \dis e^x = \sum_{n=0}^{\infty}\dfrac{1}{n!}x^n, |x|<\infty $ 
    \item $ \dis \sin x = \sum_{n=0}^{\infty}(-1)^n\dfrac{x^{2n+1}}{(2n+1)!},|x|<\infty $ 
    \item $ \dis \cos x = \sum_{n=0}^{\infty}(-1)^n\dfrac{x^{2n}}{(2n)!},|x|<\infty $ 
    \item $ \dis (1+x)^\alpha = 1 + \sum_{n=1}^{\infty}\dfrac{\mathrm{C}_\alpha^n}{n!}x^n, |x|<1 $ 
    \item $ \dis \dfrac{1}{1-x} = \sum_{n=0}^{\infty}x^n, |x|<1 $ 
    \item $ \dis \dfrac{1}{1+x} = \sum_{n=0}^{\infty}(-1)^nx^n, |x|<1 $ 
    \item $ \dis \ln(1+x) = \sum_{n=1}^{\infty} \dfrac{(-1)^{n+1}x^n}{n!}, -1<x\leq 1 $ 
    \item $ \dis \ln(1-x) = \sum_{n=1}^{\infty} \dfrac{x^n}{n!}, -1\leq x< 1 $ 
\end{itemize}

以上都展成马克劳林级数。若要展成泰勒级数，经过适当处理后可以利用马克劳林级数的结果。

展开的结果尽量合成一个级数，级数后要注收敛范围。

可以利用幂级数求数项级数，即将数项级数视为幂级数的 $ x $ 取了定值。

\section{傅里叶级数}

\subsection{傅里叶级数及傅里叶系数}

设函数 $ f(x) $ 在区间 $ [-\pi,\pi] $ 上黎曼可积，则称公式
\begin{itemize}
    \item $ \dis a_k = \dfrac{1}{\pi}\int_{-\pi}^\pi f(x)\cos kx\mathrm{d}x $ 
    \item $ \dis b_k = \dfrac{1}{\pi}\int_{-\pi}^\pi f(x)\sin kx\mathrm{d}x $ 
\end{itemize}

为函数 $ f(x) $ 的傅里叶系数。称以 $ a_k,b_k $ 为系数的三角级数$$
    \dfrac{a_0}{2} + \sum_{n=1}^{\infty}\left(a_n\cos nx + b_n\sin nx\right)
$$ 
为 $ f(x) $ 的傅里叶级数，记为
$$
    f(x)\sim \dfrac{a_0}{2} + \sum_{n=1}^{\infty}\left(a_n\cos nx + b_n\sin nx\right)
$$ 

类似地，可以定义周期为 $ 2l $ 的函数的傅里叶级数，即

\begin{itemize}
    \item $ \dis a_k = \dfrac{1}{l}\int_{-l}^l f(x)\cos \dfrac{k\pi x}{l}\mathrm{d}x $ 
    \item $ \dis b_k = \dfrac{1}{l}\int_{-l}^l f(x)\sin \dfrac{k\pi x}{l}\mathrm{d}x $
    \item $ \dis f(x) \sim \dfrac{a_0}{2} + 
    \sum_{n=1}^{\infty}\left(a_n\cos \dfrac{n\pi x}{l} + b_n\sin \dfrac{n\pi x}{l}\right) $ 
\end{itemize}

\subsection{傅里叶级数的收敛定理}

设 $ f(x) $ 是周期为 $ 2\pi $ 的周期函数，并满足迪利克雷条件，即
\begin{itemize}
    \item 在一个周期内连续或有有限个第一类间断点；
    \item 在一个周期内有有限个极值点，即单调区间个数有限，
\end{itemize}
则$ f(x) $ 的傅里叶级数收敛，且有

\begin{equation*}
    \begin{aligned}
        \dfrac{a_0}{2} + \sum_{n=1}^{\infty}\left(a_n\cos nx + b_n\sin nx\right)
        = \begin{cases}
            f(x),&x\textrm{为$ f(x) $ 的连续点}\\\dfrac{f(x-0)+f(x+0)}{2},& x\textrm{为$ f(x) $ 的间断点}
        \end{cases}
    \end{aligned}
\end{equation*}

\subsection{奇偶函数的傅里叶级数}

若 $ f(x) $ 是奇函数，则有 $ \dis a_n = \dfrac{1}{l}\int_{-l}^l f(x)\cos \dfrac{n\pi x}{l}\mathrm{d}x = 0,
f(x) = \sum_{n=1}^{\infty}b_n\sin \dfrac{n\pi x}{l} $ ，是为正弦级数。

若 $ f(x) $ 是偶函数，则有 $ \dis b_n = \dfrac{1}{l}\int_{-l}^l f(x)\sin \dfrac{n\pi x}{l}\mathrm{d}x = 0,
f(x) =\dfrac{a_0}{2}+  \sum_{n=1}^{\infty}a_n\cos \dfrac{n\pi x}{l} $ ，是为余弦级数。

以上二级数，都有 $ x\in\{f\textrm{的连续点}\} $.

\subsection{有限区间上的函数的傅里叶级数}

将有限区间上的函数展为正弦级数时，将$ f(x) $ 延拓为周期为 $ 2\pi $ 的奇函数 $ F(x) $ .
此时 $ a_n = 0,\dis b_n = \int_{-\pi}^{\pi} f(x)\sin nx \mathrm{d}x $ ，用分部积分法求 $ b_n $ .
然后有 $ \dis f(x) \xlongequal{\textrm{f(x)延拓的部分}} F(x) \xlongequal{\textrm{连续点}} \sum_{n = 1}^\infty b_n\sin nx$
并给出 $ x $ 的区间。 

展成余弦级数同理。



\end{document}
